% !PRINTED 74
\weberpage{74}
%<<CONT
\gesperrt{the index $\beta$ corresponding to a prime $q$, of all the numbers $b$, is divisible by $q$ or by $q - 1$, according as $q$ occurs several times or only once among the prime factors of $m$.}

\begin{center}
\S.\,20.

\gesperrt{The cyclotomic periods.}
\end{center}

As the simplest functions by which one may attempt to represent the cyclotomic fields, the cyclotomic periods present themselves, which are an immediate generalisation of the Gauss periods considered in \S.\,167 of the first volume. By them we understand the following.

Let $\mathfrak{A}$ denote a divisor of the group $\mathfrak{N}$ of index $e$, and let $a$ run through the numbers of $\mathfrak{A}$. If $r$ is a primitive $m$th root of unity, then the sum
\weq{1}{\eta = \overset{a}{\Sigma}\, r^a}
is called a \gesperrt{cyclotomic period belonging to the group $\mathfrak{A}$} of index $e$.

If we make in (1) one of the substitutions $(r, r^a)$, then $\eta$ remains unchanged, as follows immediately from the group property of the $a$.

In order to decompose $\mathfrak{N}$ into the subsidiary groups belonging to $\mathfrak{A}$, we must choose the numbers $1, n_1, n_2, \ldots n_{e-1}$ from $\mathfrak{N}$ suitably, so that the decomposition
\weq{2}{\mathfrak{N} = \mathfrak{A} + \mathfrak{A} n_1 + \mathfrak{A} n_2 + \cdots + \mathfrak{A} n_{e-1}}
is obtained.
If we then make in $\eta$ the substitutions
\weq{3}{(r, r), (r, r^{n_1}), \ldots (r, r^{n_{e-1}}),}
then $\eta$ passes into the conjugate periods
\weq{4}{\eta, \eta_1, \ldots \eta_{e-1}}
and if these are all different from one another, then $\eta$ belongs to the group $\mathfrak{A}$. The field $\Omega(\eta)$ belonging to $\mathfrak{A}$ consists of all rational functions of $\eta$, and the numbers $\eta_1 \ldots \eta_{e-1}$ all belong to the same group $\mathfrak{A}$.

If, however, the quantities (4) are not all different from one another, then the number $\eta$ does not belong to the group $\mathfrak{A}$, but to a more comprehensive group $\mathfrak{A}'$, of which $\mathfrak{A}$ is a divisor
%CONT>>
